\documentclass[a4paper, 11pt, reqno]{amsart}

\usepackage{amsmath, amsthm, amssymb}
\numberwithin{equation}{section}
\usepackage{xspace}
\usepackage{mathtools}
\usepackage{mathrsfs}
\usepackage{hyperref}

\theoremstyle{plain}
\newtheorem{thm}{Theorem}[section]
\newtheorem{lemma}[thm]{Lemma}
\newtheorem{corollary}[thm]{Corollary}
\newtheorem{prop}[thm]{Proposition}
\theoremstyle{remark}
\newtheorem*{claim}{Claim}
\theoremstyle{definition}
\newtheorem{remark}[thm]{Remark}
\newtheorem{defn}[thm]{Definition}

\newcommand{\C}{\mathbb{C}}
\newcommand{\B}{\mathbb{B}}
\newcommand{\KKK}{{\mathbb K}}
\newcommand{\KK}{{\mathcal K}}
\newcommand{\XX}{{\mathbb{X}}}
\newcommand{\VV}{{\mathcal V}}
\newcommand{\EE}{{\mathcal E}}
\newcommand{\GG}{{\mathcal G}}
\newcommand{\S}{\mathbb{S}}
\let\restr\relax
\newcommand{\la}{\langle}
\newcommand{\ra}{\rangle}
\DeclareMathOperator{\Ad}{Ad}

\begin{document}

\title[Properly proximal graph product von Neumann algebras]{Properly proximal graph product von Neumann algebras}
\maketitle

\section{Introduction}

The notion of proper proximality for groups was introduced in \cite{BIP21}, generalizing Ozawa's notion of biexactness from \cite{Oza04}. This notion is a strengthening of non-inner-amenability in the sense of Effros, and has applications to the study of von Neumann algebras. In particular, it is an invariant up to measure equivalence and von Neumann equivalence \cite{IsPeRu19}, and when combined with W$^*$CBAP it yields absence of Cartan subalgebras, following ideas of \cite{OzPo10a}.

In the paper \cite{DKE24properproximalgroups}, proper proximality was classified for graph products of groups, in the sense of Green \cite{Gr90}. In this note this is generalized to the setting of graph products of tracial von Neumann algebras \cite{CaFi17}. Indeed the notion of proper proximality has been extended to von Neumann algebras in \cite{DKEP23}, so this question is natural. We prove the following result.

\begin{thm}\label{thm:main}
Let $\GG$ be a finite graph with $|\VV(\GG)|\geq 3$, and let $(M_v,\varphi_v)$ be a separable tracial von Neumann algebra for each $v\in\VV(\GG)$. If $\GG$ is irreducible (in the sense that there do not exist subgraphs $\GG_1,\GG_2$ with $\VV(\GG_1)\cup\VV(\GG_2)=\VV(\GG)$ such that $(v,u)\in\EE(\GG)$ for all $v\in\VV(\GG_1)$ and $u\in\VV(\GG_2)$ with $v\neq u$), and each $M_v$ has a trace-zero unitary, then $M=\bigstar_{\GG}M_v$ is properly proximal.
\end{thm}

A full classification for graph products then follows by combining this statement with the known fact from \cite{DKEP23} that a tensor product of finite von Neumann algebras $N_1\overline{\otimes}N_2$ is properly proximal if and only if both are: every finite graph decomposes into joins (corresponding to tensor products) of irreducible subgraphs, and each irreducible factor with at least three vertices is handled by Theorem~\ref{thm:main}.

For the sake of brevity we follow the notation of \cite{DKEP23}, and assume familiarity with the standard notions and arguments of von Neumann algebra theory. We recall in particular the boundary algebra $\widetilde{\S}_{\XX}(M)$ relative to an $M$-boundary piece $\XX$, the bidual $\B(L^2M)^{\sharp*}_J$, the canonical normal embedding $\iota^\sharp\colon\B(L^2M)\to\B(L^2M)^{\sharp*}_J$, and the characterisation of proper proximality relative to $\XX$ via the absence of an $M$-central state on $\widetilde\S_{\XX}(M)$ whose restriction to $M$ is normal \cite[Lemma 8.5]{DKEP23}.

\section{Commutation of support projections}\label{sec:supports}

Let $(\GG,\EE)$ be a finite graph and let $\{(M_v,\varphi_v):v\in\VV(\GG)\}$ be a family of von Neumann algebras with faithful normal states. We write $(M,\varphi)$ for the graph product von Neumann algebra, and for a subgraph $\GG'\subseteq\GG$ we write $\Gamma(\GG')=\bigstar_{\GG'}M_v\subseteq M$ for the corresponding subalgebra; we always identify $L^2\Gamma(\GG')$ with the closed subspace of $L^2M$ spanned by the reduced words supported on $\VV(\GG')$.

Fix two subgraphs $\GG_1,\GG_2$, set $M_i=\Gamma(\GG_i)$ for $i=1,2$, and let $N=\Gamma(\GG_1\cap\GG_2)=M_1\cap M_2$. Let $e_i$ be the orthogonal projection of $L^2M$ onto $L^2M_i$, and let $e$ be the orthogonal projection of $L^2M$ onto $L^2N$. We record the \emph{commuting-square identity}
\begin{equation}\label{eq:commsquare}
e_1e_2=e_2e_1=e,
\end{equation}
which holds because a reduced word $\xi_{w_1\cdots w_k}$ survives $e_2$ iff all its letters lie in $\VV(\GG_2)$ and then survives $e_1$ iff all its letters also lie in $\VV(\GG_1)$, i.e.\ iff all letters lie in $\VV(\GG_1\cap\GG_2)$; this is exactly the range of $e$ (equivalently, $E_{M_1}E_{M_2}=E_{M_2}E_{M_1}=E_N$, the standard commuting square for graph products, cf.\ \cite[\S3]{CaFi17}).

Consider the $M$-boundary pieces
\[
\XX_{M_i}=\overline{MJMJ\,e_i\,\B(L^2M)\,e_i\,JMJM}^{\,\|\cdot\|},\qquad
\XX_{N}=\overline{MJMJ\,e\,\B(L^2M)\,e\,JMJM}^{\,\|\cdot\|},
\]
together with their $M$-$M$, $JMJ$-$JMJ$ closures $\KKK^{\infty,1}_{M_i}(M)$ and $\KKK^{\infty,1}_{N}(M)$. The map $\iota^\sharp\colon\B(L^2M)\to\B(L^2M)^{\sharp*}_J$ is the canonical embedding. Let
\[
q_i=\bigvee_{u,w\in\mathcal U(M)}\iota^\sharp(uJwJ\,e_i\,Jw^*Ju^*),\qquad
q=\bigvee_{u,w\in\mathcal U(M)}\iota^\sharp(uJwJ\,e\,Jw^*Ju^*)
\]
be the identities (central support projections) of $(\XX_{M_i})^{\sharp*}_J$ and $(\XX_N)^{\sharp*}_J$ inside $\B(L^2M)^{\sharp*}_J$.

Our aim in this section is to prove that $q_1q_2=q_2q_1=q$. Since $\iota^\sharp$ is a normal $*$-homomorphism, we may and do realise $\B(L^2M)^{\sharp*}_J$ on a Hilbert space $\KK$ via a faithful normal representation; for $a\in\B(L^2M)^{\sharp*}_J$ and a subset $S\subseteq\B(L^2M)$ we abbreviate $\iota^\sharp(S)\KK=\overline{\operatorname{span}}\{\iota^\sharp(x)\xi:x\in S,\ \xi\in\KK\}$. Because $e\in e\B(L^2M)e\subseteq\XX_N$ and $q$ is the identity of $(\XX_N)^{\sharp*}_J=q\B(L^2M)^{\sharp*}_J q$, we have
\begin{equation}\label{eq:qrange}
q\KK=\overline{\iota^\sharp(MJMJe)\KK},\qquad q\,\iota^\sharp(e)=\iota^\sharp(e)=\iota^\sharp(e)\,q .
\end{equation}
Indeed, $\iota^\sharp(uJwJeJw^*Ju^*)\leq q$ has range $uJwJ\,\iota^\sharp(e)\KK$, and taking the join over $u,w\in\mathcal U(M)$ and using that $M$ (resp.\ $JMJ$) is the closed linear span of its unitaries gives the first identity in \eqref{eq:qrange}; the second follows because $\iota^\sharp(e)\leq q$.

\begin{lemma}\label{lem:range}
With the above representation $\KK$,
\[
\iota^\sharp(e_1)\,q_2\,\KK=\overline{\iota^\sharp(e_1MJMJe_2)\KK}\subseteq\overline{\iota^\sharp(MJMJe)\KK}=q\KK .
\]
\end{lemma}

\begin{proof}
The first equality holds because $q_2\KK=\overline{\iota^\sharp(MJMJe_2)\KK}$ (the analogue of \eqref{eq:qrange} for $\GG_2$), and because $\iota^\sharp$ is multiplicative. The last equality is \eqref{eq:qrange}. It remains to prove the inclusion, for which it suffices to show that
\begin{equation}\label{eq:goal}
\iota^\sharp(e_1MJMJe_2)\KK\subseteq\overline{\iota^\sharp(MJMJe)\KK}=q\KK .
\end{equation}

We first decompose $e_1MJMJe_2$. Since $M=M_1\oplus(M\ominus M_1)$ orthogonally with $e_1$ and $Je_1J$ the corresponding $L^2$-projections, and since $e_1$ commutes with $M_1$ and with $JM_1J$, we have
\begin{align}
e_1MJMJe_2
&=e_1(M\ominus M_1)J(M\ominus M_1)Je_2
 +e_1M_1J(M\ominus M_1)Je_2
 +e_1MJM_1Je_2 \nonumber\\
&=e_1(M\ominus M_1)J(M\ominus M_1)Je_2
 +M_1e_1J(M\ominus M_1)Je_2 \nonumber\\
&\quad
 +JM_1J\,e_1(M\ominus M_1)e_2
 +JM_1J\,M_1\,e .\label{eq:decomp}
\end{align}
Here the first equality is the orthogonal expansion of $MJMJ=\big(M_1+(M\ominus M_1)\big)J\big(M_1+(M\ominus M_1)\big)J$, after collecting the two terms containing $JM_1J$ into $MJM_1J$. For the second equality we used $e_1M_1=M_1e_1$ in the second summand, and in the third we split $e_1MJM_1Je_2=JM_1J\,e_1MJ\!\cdot\!1\,e_2$ as
\[
e_1MJM_1Je_2=JM_1J\,e_1MJ\!\cdot\!1\,e_2=JM_1J\big(e_1M_1e_2+e_1(M\ominus M_1)e_2\big)
=JM_1J\,M_1\,e+JM_1J\,e_1(M\ominus M_1)e_2,
\]
where we used that $JM_1J$ commutes with $e_1,e_2$ and with $M\ominus M_1$, the orthogonal split $M=M_1\oplus(M\ominus M_1)$, and finally the identity $e_1M_1e_2=M_1\,e_1e_2=M_1\,e$, which follows from $e_1M_1=M_1e_1$ together with the commuting-square identity \eqref{eq:commsquare}. We now justify, term by term, that the range of each summand in \eqref{eq:decomp} lies in $q\KK$.

\medskip\noindent\textbf{Term 1.}
We claim that
\begin{equation}\label{eq:orth}
e_2\,(M\ominus M_1)J(M\ominus M_1)J\,(e_1-e)=0,
\end{equation}
equivalently, by taking adjoints, $(e_1-e)(M\ominus M_1)J(M\ominus M_1)J\,e_2=0$. Recall that $(e_1-e)L^2M=L^2M_1\ominus L^2N$ is the closed span of the reduced words $\xi_{w_1\cdots w_k}$ all of whose letters lie in $\VV(\GG_1)$, with at least one letter in $\VV(\GG_1)\setminus\VV(\GG_1\cap\GG_2)=\VV(\GG_1)\setminus\VV(\GG_2)$. For $a,b\in M\ominus M_1$ (left action of $a$, right action of $JbJ$) and such a word $\xi$, the vector $aJbJ\xi$ is, by the reduced-word multiplication rules in a graph product \cite[\S2]{CaFi17}, a linear combination of reduced words each of which contains at least one letter that is \emph{not} in $\VV(\GG_2)$: this letter is supplied by the original letter of $\xi$ lying in $\VV(\GG_1)\setminus\VV(\GG_2)$, which cannot be cancelled or absorbed by letters coming from $a,b\in M\ominus M_1$ (these introduce only further letters outside $\VV(\GG_1)$, hence outside $\VV(\GG_1\cap\GG_2)$, and cancellation against the distinguished letter would require a letter in the same vertex group, impossible since $a,b\notin M_1$). Any reduced word containing a letter outside $\VV(\GG_2)$ is orthogonal to $L^2M_2=e_2L^2M$, which proves \eqref{eq:orth}.

Consequently $e_1(M\ominus M_1)J(M\ominus M_1)Je_2=e(M\ominus M_1)J(M\ominus M_1)Je_2$, because $e_1=e+(e_1-e)$ and the $(e_1-e)$-part is killed by \eqref{eq:orth} after pairing with $e_2$ on the right. Since $e\in\XX_N$ and, by \eqref{eq:qrange}, $q\iota^\sharp(e)=\iota^\sharp(e)$, the range of $\iota^\sharp\big(e(M\ominus M_1)J(M\ominus M_1)Je_2\big)$ is contained in $\iota^\sharp(e)\KK\subseteq q\KK$. Thus Term~1 has range in $q\KK$.

\medskip\noindent\textbf{Term 2.}
By the same reduced-word analysis as for \eqref{eq:orth} (now with $a\in M\ominus M_1$ acting on the right via $J(M\ominus M_1)J$ and the identity on the left), we have $(e_1-e)J(M\ominus M_1)Je_2=0$, hence $e_1J(M\ominus M_1)Je_2=eJ(M\ominus M_1)Je_2$. Therefore
\[
M_1e_1J(M\ominus M_1)Je_2=M_1\,eJ(M\ominus M_1)Je_2 .
\]
Now $eJ(M\ominus M_1)Je_2\in e\,\B(L^2M)$ has its left support under $e$, so it lies in $e\B(L^2M)e_2\subseteq\overline{MJMJ\,e\,\B(L^2M)\,e_2\,JMJM}$. Multiplying on the left by $M_1\subseteq M$ keeps the operator in $\XX_N$'s closure of $MJMJ\,e\,\B(L^2M)\,e\,JMJM$ after using $e_2=e_2 e_2$ and $e_2\in\B(L^2M)$; more directly, $q\,\iota^\sharp(M_1\,e\,X)=\iota^\sharp(M_1)\,q\,\iota^\sharp(e)\,\iota^\sharp(X)=\iota^\sharp(M_1\,e\,X)$ for any $X\in\B(L^2M)$, since $q$ commutes with $\iota^\sharp(M)$ and $q\iota^\sharp(e)=\iota^\sharp(e)$. Hence the range of $\iota^\sharp(M_1e_1J(M\ominus M_1)Je_2)$ lies in $q\KK$.

\medskip\noindent\textbf{Term 3.}
By the same argument with sides exchanged, $e_1(M\ominus M_1)(e_2-e\,?)$\,: precisely, the reduced-word analysis gives $(e_1-e)(M\ominus M_1)e_2=0$, hence $e_1(M\ominus M_1)e_2=e(M\ominus M_1)e_2$. Therefore
\[
JM_1J\,e_1(M\ominus M_1)e_2=JM_1J\,e(M\ominus M_1)e_2 .
\]
Since $q$ commutes with $\iota^\sharp(JMJ)\supseteq\iota^\sharp(JM_1J)$ and $q\iota^\sharp(e)=\iota^\sharp(e)$, the same computation as in Term~2 shows
\[
q\,\iota^\sharp\big(JM_1J\,e(M\ominus M_1)e_2\big)=\iota^\sharp\big(JM_1J\,e(M\ominus M_1)e_2\big),
\]
so this summand has range in $q\KK$.

\medskip\noindent\textbf{Term 4.}
The summand $JM_1J\,M_1\,e$ already carries the factor $e$ on the right. As $q$ commutes with $\iota^\sharp(M)$ and $\iota^\sharp(JMJ)$ and $q\iota^\sharp(e)=\iota^\sharp(e)$,
\[
q\,\iota^\sharp(JM_1J\,M_1\,e)=\iota^\sharp(JM_1J\,M_1)\,q\,\iota^\sharp(e)=\iota^\sharp(JM_1J\,M_1\,e),
\]
so its range lies in $q\KK$.

\medskip
Adding the four contributions and using \eqref{eq:decomp}, we obtain \eqref{eq:goal}, which completes the proof.
\end{proof}

\begin{corollary}\label{cor:commute}
$[q_1,q_2]=0$ and $q_1q_2=q_2q_1=q$.
\end{corollary}

\begin{proof}
Since $q_2$ commutes with $\iota^\sharp(M)$ and $\iota^\sharp(JMJ)$ (it is, by construction, in the relative commutant of these algebras, being the central support of an $M$-$M$, $JMJ$-$JMJ$ boundary piece), and since
$q_1=\bigvee_{u,w\in\mathcal U(M)}\iota^\sharp(uJwJ\,e_1\,Jw^*Ju^*)$,
to prove $[q_1,q_2]=0$ it suffices to show $[\iota^\sharp(e_1),q_2]=0$.

By Lemma~\ref{lem:range}, $\iota^\sharp(e_1)q_2\KK\subseteq q\KK$, so $q\,\iota^\sharp(e_1)q_2=\iota^\sharp(e_1)q_2$. As $q\le q_2$ (because $\XX_N\subseteq\XX_{M_2}$, since $e=e e_2=e_2 e$ shows $e\in e_2\B(L^2M)e_2$ and hence $\XX_N\subseteq\XX_{M_2}$, giving $q\le q_2$), we get
\[
\iota^\sharp(e_1)q_2=q\,\iota^\sharp(e_1)q_2=q_2\,q\,\iota^\sharp(e_1)q_2 ,
\]
where the last equality uses $q\le q_2$ so $q_2q=q$. In particular $\iota^\sharp(e_1)q_2=q_2\,\big(q\,\iota^\sharp(e_1)q_2\big)=q_2\,\iota^\sharp(e_1)q_2$. Taking adjoints (recall $e_1=e_1^*$, $q_2=q_2^*$) gives $q_2\,\iota^\sharp(e_1)=q_2\,\iota^\sharp(e_1)q_2=\iota^\sharp(e_1)q_2$, i.e.\ $[\iota^\sharp(e_1),q_2]=0$. Therefore $[q_1,q_2]=0$.

We now compute $q_1q_2$. Using $[\iota^\sharp(e_1),q_2]=0$ and the established identities $\iota^\sharp(e_1)q_2=q\,\iota^\sharp(e_1)q_2=q\,\iota^\sharp(e_1)q$ (the last because $q\le q_2$ gives $q_2q=q$, so $q\iota^\sharp(e_1)q_2=q\iota^\sharp(e_1)q_2 q\cdot$\,; more simply $q\iota^\sharp(e_1)q_2$ multiplied on the right by $q\le q_2$ is unchanged after inserting $q_2q=q$), we get for all $u,w\in\mathcal U(M)$,
\[
\iota^\sharp(uJwJ\,e_1\,Jw^*Ju^*)\,q_2
=\iota^\sharp(uJwJ)\,\big(\iota^\sharp(e_1)q_2\big)\,\iota^\sharp(Jw^*Ju^*)
=\iota^\sharp(uJwJ)\,q\,\iota^\sharp(e_1)q\,\iota^\sharp(Jw^*Ju^*).
\]
Since $q$ commutes with $\iota^\sharp(M)$ and $\iota^\sharp(JMJ)$,
\[
q_1q_2=\bigvee_{u,w}\iota^\sharp(uJwJ\,e_1\,Jw^*Ju^*)\,q_2
=q\Big(\bigvee_{u,w}\iota^\sharp(uJwJ\,e_1\,Jw^*Ju^*)\Big)q
=q\,q_1\,q=q,
\]
the last equality because $q\le q_1$ (as $\XX_N\subseteq\XX_{M_1}$ for the same reason as $\XX_N\subseteq\XX_{M_2}$, using $e=ee_1=e_1e$). By symmetry $q_2q_1=q$, and we have already shown the two commute.
\end{proof}

\begin{corollary}\label{cor:intersection}
For any two subgraphs $\GG_1,\GG_2\subseteq\GG$,
\[
\KKK^{1,\infty}_{M_1}(M)\cap\KKK^{1,\infty}_{M_2}(M)=\KKK^{1,\infty}_{N}(M),
\qquad N=\Gamma(\GG_1\cap\GG_2).
\]
Moreover, for any finite family $\GG_1,\dots,\GG_s\subseteq\GG$ the support projections commute pairwise and
\begin{equation}\label{eq:meet}
\bigwedge_{i=1}^s q_{\XX_{\Gamma(\GG_i)}}=q_{\XX_{\Gamma(\cap_i\GG_i)}},\qquad\text{equivalently}\qquad
\Big(\bigwedge_{i=1}^s q_{\XX_{\Gamma(\GG_i)}}\Big)^{\!\perp}=\bigvee_{i=1}^s q_{\XX_{\Gamma(\GG_i)}}^{\perp}.
\end{equation}
\end{corollary}

\begin{proof}
Since $\KKK^{1,\infty}_{M_i}(M)=(\iota^\sharp)^{-1}\big((\XX_{M_i})^{\sharp*}_J\big)$ and $(\XX_{M_i})^{\sharp*}_J=q_i\,\B(L^2M)^{\sharp*}_J\,q_i$, the equality $\KKK^{1,\infty}_{M_1}(M)\cap\KKK^{1,\infty}_{M_2}(M)=\KKK^{1,\infty}_{N}(M)$ is equivalent to $q_1\B(L^2M)^{\sharp*}_Jq_1\cap q_2\B(L^2M)^{\sharp*}_Jq_2=q\B(L^2M)^{\sharp*}_Jq$. By Corollary~\ref{cor:commute} the central projections $q_1,q_2$ commute and $q_1q_2=q$, so $q_1\B(L^2M)^{\sharp*}_Jq_1\cap q_2\B(L^2M)^{\sharp*}_Jq_2=(q_1\wedge q_2)\B(L^2M)^{\sharp*}_J(q_1\wedge q_2)=q\B(L^2M)^{\sharp*}_Jq$, as desired.

For the family statement we argue by induction on $s$. The case $s=1$ is trivial and the case $s=2$ is Corollary~\ref{cor:commute} (the meet of two commuting projections is their product). Assume \eqref{eq:meet} for $s-1$ subgraphs and set $\GG'=\bigcap_{i=2}^s\GG_i$, $M'=\Gamma(\GG')$. By the induction hypothesis $\bigwedge_{i=2}^s q_{\XX_{\Gamma(\GG_i)}}=q_{\XX_{M'}}$ and these projections commute pairwise. Applying Corollary~\ref{cor:commute} to the two subgraphs $\GG_1$ and $\GG'$ gives $[q_{\XX_{M_1}},q_{\XX_{M'}}]=0$ and $q_{\XX_{M_1}}\wedge q_{\XX_{M'}}=q_{\XX_{\Gamma(\GG_1\cap\GG')}}=q_{\XX_{\Gamma(\cap_{i=1}^s\GG_i)}}$. Hence $\bigwedge_{i=1}^s q_{\XX_{\Gamma(\GG_i)}}=q_{\XX_{M_1}}\wedge q_{\XX_{M'}}=q_{\XX_{\Gamma(\cap_i\GG_i)}}$, and $q_{\XX_{M_1}}$ commutes with each of $q_{\XX_{\Gamma(\GG_2)}},\dots,q_{\XX_{\Gamma(\GG_s)}}$ by Corollary~\ref{cor:commute}, so the whole family commutes pairwise. The equivalent De Morgan form follows by taking orthocomplements of commuting projections.
\end{proof}

\section{Stability of relative proper proximality under intersections}

\begin{lemma}\label{lem:intersect-pp}
Suppose $(M_v,\varphi_v)$ is separable and tracial for each $v\in\VV(\GG)$. Let $\{\GG_i\}_{i=1}^s$ be subgraphs of $\GG$. If $M$ is properly proximal relative to $\Gamma(\GG_i)$ for each $i$, then $M$ is properly proximal relative to $\Gamma\big(\bigcap_i\GG_i\big)$.
\end{lemma}

\begin{proof}
Write $\XX_i=\XX_{\Gamma(\GG_i)}$ and $\XX_\cap=\XX_{\Gamma(\cap_i\GG_i)}$, with support projections $q_i:=q_{\XX_i}$ and $q_\cap:=q_{\XX_\cap}$ in $\B(L^2M)^{\sharp*}_J$. By Corollary~\ref{cor:intersection} these commute pairwise and
\begin{equation}\label{eq:demorgan}
q_\cap=\bigwedge_i q_i,\qquad q_\cap^\perp=\bigvee_i q_i^\perp .
\end{equation}

Suppose, for contradiction, that $M$ is \emph{not} properly proximal relative to $\Gamma(\cap_i\GG_i)$. By \cite[Lemma 8.5]{DKEP23} there is an $M$-central state $\psi$ on $\widetilde\S_{\XX_\cap}(M)$ whose restriction to $M$ is normal (equal to the trace $\tau$). Recall $\widetilde\S_{\XX_\cap}(M)\subseteq\B(L^2M)^{\sharp*}_J$ and that conjugation by $\iota^\sharp(\mathcal U(M))$ leaves $\widetilde\S_{\XX_\cap}(M)$ invariant.

\emph{Case 1: $\psi(q_\cap)\neq0$.} Then $\psi_0:=\psi(q_\cap)^{-1}\,\psi\circ\Ad(q_\cap)$ is a state. Since $q_\cap$ is the central support of $\XX_\cap$ and commutes with $\iota^\sharp(M)$ and $\iota^\sharp(JMJ)$, $\Ad(q_\cap)$ maps $\B(L^2M)^{\sharp*}_J$ into $(\XX_\cap)^{\sharp*}_J\subseteq(\XX_i)^{\sharp*}_J\subseteq\widetilde\S_{\XX_i}(M)$ (the inclusion $(\XX_\cap)^{\sharp*}_J\subseteq(\XX_i)^{\sharp*}_J$ is $q_\cap\le q_i$ from \eqref{eq:demorgan}). Moreover $\psi_0$ is $M$-central because $\psi$ is $M$-central and $q_\cap$ commutes with $\iota^\sharp(M)$; and $\psi_0|_M$ is a positive normal multiple of $\tau$, hence after the normalisation it equals $\tau$ on $M$ (using $\Ad(q_\cap)\iota^\sharp(x)=q_\cap\iota^\sharp(x)$ and $M$-centrality: $\psi(q_\cap\iota^\sharp(x))=\psi(\iota^\sharp(x)q_\cap)$, so $\psi_0(\iota^\sharp(x))=\psi(q_\cap)^{-1}\psi(\iota^\sharp(x)q_\cap)$, and $x\mapsto\psi(\iota^\sharp(x)q_\cap)$ is normal on $M$). Thus $\psi_0$ is an $M$-central state on $\widetilde\S_{\XX_i}(M)$, normal on $M$, contradicting proper proximality of $M$ relative to $\Gamma(\GG_i)$.

\emph{Case 2: $\psi(q_\cap)=0$, i.e.\ $\psi(q_\cap^\perp)=1$.} By \eqref{eq:demorgan}, $q_\cap^\perp=\bigvee_i q_i^\perp$. Since $\psi$ is a state and $q_\cap^\perp$ is a join of finitely many commuting projections $q_i^\perp$, we have $q_\cap^\perp\le\sum_i q_i^\perp$, so $1=\psi(q_\cap^\perp)\le\sum_i\psi(q_i^\perp)$, and there exists $i_0$ with $\psi(q_{i_0}^\perp)\neq0$. Set $\psi_1:=\psi(q_{i_0}^\perp)^{-1}\,\psi\circ\Ad(q_{i_0}^\perp)$.

We claim $\Ad(q_{i_0}^\perp)$ maps $\widetilde\S_{\XX_\cap}(M)$ into $\widetilde\S_{\XX_{i_0}}(M)$. Indeed $q_{i_0}^\perp$ commutes with $\iota^\sharp(JMJ)$ (since $q_{i_0}$ does), so
\[
\Ad(q_{i_0}^\perp)\big(\widetilde\S_{\XX_\cap}(M)\big)\subseteq q_{i_0}^\perp\,\B(L^2M)^{\sharp*}_J\,q_{i_0}^\perp\cap\iota^\sharp(JMJ)' .
\]
By the definition of $\widetilde\S_{\XX_{i_0}}(M)$ as the elements of $\B(L^2M)^{\sharp*}_J\cap\iota^\sharp(JMJ)'$ supported on $q_{i_0}^\perp$ together with $(\XX_{i_0})^{\sharp*}_J$ \cite[\S8]{DKEP23}, the right-hand side is contained in $\widetilde\S_{\XX_{i_0}}(M)$. Hence $\psi_1$ is a state on $\widetilde\S_{\XX_{i_0}}(M)$. It is $M$-central because $q_{i_0}^\perp$ commutes with $\iota^\sharp(M)$ and $\psi$ is $M$-central, and $\psi_1|_M=\tau$ by the same normalisation argument as in Case~1. This contradicts proper proximality of $M$ relative to $\Gamma(\GG_{i_0})$.

In both cases we reach a contradiction, so $M$ is properly proximal relative to $\Gamma(\cap_i\GG_i)$.
\end{proof}

\section{The amalgamated free product step}

\begin{lemma}\label{lem:afp}
Let $M_1,M_2$ be separable tracial von Neumann algebras with a common von Neumann subalgebra $B$ that is the range of trace-preserving conditional expectations $E_B^{M_i}$. Suppose there exist unitaries $u_1,u_2\in\mathcal U(M_1)$ and $u_3\in\mathcal U(M_2)$ with
\begin{equation}\label{eq:eb0}
E_B(u_1)=E_B(u_2)=E_B(u_3)=E_B(u_1^*u_2)=0 .
\end{equation}
Then the amalgamated free product $M:=M_1*_BM_2$ is properly proximal relative to $B$.
\end{lemma}

\begin{proof}
We use the projections $\lambda_i(e_B)$ defined in \cite{toyosawa2025relative}: $\lambda_i(e_B)$ is the orthogonal projection of $L^2M$ onto the closed span of the reduced words that begin (on the left) with a letter from $L^2M_{i+1}^o$ (indices mod $2$, so $M_3:=M_1$) or that lie in $L^2B$. By \cite[Prop.~4.16]{toyosawa2025relative}, $\lambda_i(e_B)\in\S_B(M)$, and in particular $\iota^\sharp(\lambda_i(e_B))\in\widetilde\S_B(M)$.

Suppose, for contradiction, that $M$ is not properly proximal relative to $B$. By \cite[Lemma 8.5]{DKEP23} there is an $M$-central state $\varphi$ on $\widetilde\S_B(M)$ whose restriction to $M$ is the (normal) trace.

\medskip\noindent\textbf{Step 1: $\varphi(\iota^\sharp(e_B))=0$.}
Suppose not. The projections $p_u:=\iota^\sharp(u e_B u^*)$ for $u\in\mathcal U(M)$ satisfy, by $M$-centrality of $\varphi$, $\varphi(p_u)=\varphi(\iota^\sharp(e_B))>0$. Let $p:=\bigvee_{u\in\mathcal U(M)}p_u=\iota^\sharp\big(\bigvee_{u}ue_Bu^*\big)$, a projection commuting with $\iota^\sharp(M)$. Then $\varphi\circ\Ad(p)$ is an $M$-central positive functional on $\B(L^2M)^{\sharp*}_J$, normal on $M$, that does not vanish on $\iota^\sharp(Me_BM)$ (since $\varphi(p)\ge\varphi(\iota^\sharp(e_B))>0$). Pulling back along $\iota^\sharp$, this produces a nonzero $M$-central positive functional on $\B(L^2M)$ that does not vanish on $Me_BM$. By \cite[Lemma 3.3]{OP10b} (equivalently \cite[Theorem 2]{BH18amenableabsorption}), this forces $M\preceq_M B$ in the sense of Popa's intertwining. But $M\not\preceq_M B$: for the unitary $w:=u_1u_2\in\mathcal U(M_1)$ one checks $\|E_B(x\,w^n\,y)\|_2\to0$ as $n\to\infty$ for all $x,y\in M$. Indeed $E_B(u_1u_2)=E_B(u_1)E_B(u_2)+\text{(words of positive length)}$ has the reduced form $u_1^o u_2^o$ with $u_i^o=u_i-E_B(u_i)=u_i$ by \eqref{eq:eb0}, so $w=u_1u_2$ is a single reduced word alternating between $M_1^o$ and $M_2^o$ of length $2$; hence $w^n$ is reduced of length $2n$, and the $L^2$-norm of $B$-components of $xw^ny$ tends to $0$ by the asymptotic orthogonality of long reduced words in an amalgamated free product (this is the standard mixing criterion, cf.\ \cite[\S2]{toyosawa2025relative}). This contradiction shows $\varphi(\iota^\sharp(e_B))=0$.

Here we used that $u_1,u_2\in M_1^o$ (as $E_B(u_1)=E_B(u_2)=0$) and that $u_1u_2$ stays reduced because $E_B(u_1^*u_2)=0$ guarantees $u_1^*u_2\in M_1^o$, so no cancellation collapses $w^n$; concretely $w^n=(u_1u_2)^n$ written as $u_1\,(u_2u_1)\cdots(u_2u_1)\,u_2$ is an alternating reduced word of length $2n$ because each $u_2u_1$ and $u_1u_2$ has nonzero $M_1^o$/$M_2^o$ letters, none of which lies in $B$.

\medskip\noindent\textbf{Step 2: the tree-like inequalities.}
Every nonzero reduced word in $M$ either lies in $L^2B$ or begins with a letter in $L^2M_1^o$ or in $L^2M_2^o$. Counting multiplicities ($L^2B$ is counted by both $\lambda_1(e_B)$ and $\lambda_2(e_B)$ since each includes the summand $L^2B$), this gives the operator identity
\begin{equation}\label{eq:lam-sum}
\lambda_1(e_B)+\lambda_2(e_B)=1+e_B .
\end{equation}
Moreover, the action of a $B$-orthogonal unitary moves the "beginning letter'' of a reduced word, yielding (cf.\ \cite[\S4]{toyosawa2025relative})
\begin{equation}\label{eq:tree}
u_1^*\,\lambda_1(e_B)\,u_1+u_2^*\,\lambda_1(e_B)\,u_2\le\lambda_2(e_B),\qquad
u_3^*\,\lambda_2(e_B)\,u_3\le\lambda_1(e_B).
\end{equation}
We justify \eqref{eq:tree}. The projection $\lambda_1(e_B)$ projects onto words beginning in $M_2^o$ or in $B$. For $u_1\in M_1^o$, a word $\xi$ beginning in $M_2^o$ or in $B$ satisfies that $u_1\xi$ begins in $M_1^o$ (the new leftmost letter $u_1\in M_1^o$ does not cancel, since $\xi$ begins in $M_2^o$ or $B$); hence $u_1\xi\in\operatorname{ran}\lambda_2(e_B)$. Thus $u_1\lambda_1(e_B)u_1^*\le\lambda_2(e_B)$ — equivalently $u_1^*\lambda_2(e_B)u_1\ge\lambda_1(e_B)$; rewriting with $\lambda_1$ on the inside and using that $u_1,u_2$ map the $\lambda_1$-space into mutually orthogonal subspaces of the $\lambda_2$-space (because $u_1,u_2\in M_1^o$ with $E_B(u_1^*u_2)=0$ makes $u_1\eta\perp u_2\eta'$ for $\eta,\eta'$ in the relevant subspaces), we obtain the first inequality of \eqref{eq:tree}. The second inequality is the same statement with the roles of the two free factors exchanged, using $u_3\in M_2^o$.

\medskip\noindent\textbf{Step 3: contradiction.}
Applying the state $\varphi$ to \eqref{eq:tree} and using $M$-centrality (so $\varphi(u^*\,\iota^\sharp(\lambda_i(e_B))\,u)=\varphi(\iota^\sharp(\lambda_i(e_B)))$ for any $u\in\mathcal U(M)$) and positivity of $\varphi$, write $a:=\varphi(\iota^\sharp(\lambda_1(e_B)))$ and $b:=\varphi(\iota^\sharp(\lambda_2(e_B)))$. The first inequality of \eqref{eq:tree} gives $2a\le b$; the second gives $b\le a$. Hence $2a\le b\le a$, and since $a,b\ge0$ this forces $a=b=0$. On the other hand, applying $\varphi$ to \eqref{eq:lam-sum} and using Step~1 ($\varphi(\iota^\sharp(e_B))=0$) and $\varphi(1)=1$,
\[
a+b=\varphi\big(\iota^\sharp(\lambda_1(e_B)+\lambda_2(e_B))\big)=\varphi(\iota^\sharp(1)+\iota^\sharp(e_B))=1+0=1,
\]
contradicting $a=b=0$. Therefore $M$ is properly proximal relative to $B$.
\end{proof}

\section{Proof of the main theorem}

\begin{proof}[Proof of Theorem~\ref{thm:main}]
For a vertex $v\in\VV(\GG)$ let $\operatorname{link}(v)=\{u\in\VV(\GG):u\neq v,\ (u,v)\in\EE(\GG)\}$ and $\operatorname{star}(v)=\{v\}\cup\operatorname{link}(v)$. The basic graph-product amalgamation \cite{Gr90,CaFi17} gives, for each $v$,
\begin{equation}\label{eq:star-decomp}
M=\Gamma(\GG\setminus\{v\})*_{\Gamma(\operatorname{link}(v))}\Gamma(\operatorname{star}(v)),
\end{equation}
where $\GG\setminus\{v\}$ denotes the full subgraph on $\VV(\GG)\setminus\{v\}$. Put $M_1=\Gamma(\GG\setminus\{v\})$, $B=\Gamma(\operatorname{link}(v))$ and $M_2=\Gamma(\operatorname{star}(v))$; note $B\subseteq M_1$ and $B\subseteq M_2$, both as ranges of trace-preserving expectations.

\medskip\noindent\textbf{The factor $M_2$ provides $u_3$.}
Since $v\in\operatorname{star}(v)$ but $v\notin\operatorname{link}(v)$, the vertex algebra $M_v\subseteq M_2$ is not contained in $B$, and a trace-zero unitary $a_v\in M_v$ (which exists by hypothesis) satisfies $E_B(a_v)=\varphi(a_v)\,1=0$: any single letter from $M_v^o$ with $v\notin\operatorname{link}(v)$ is orthogonal to $L^2B$. Set $u_3:=a_v$.

\medskip\noindent\textbf{The factor $M_1$ provides $u_1,u_2$.}
Let $W:=\VV(\GG)\setminus\operatorname{star}(v)=\{w\neq v:(w,v)\notin\EE(\GG)\}$ be the set of vertices distinct from $v$ and not adjacent to $v$; note $W\subseteq\VV(\GG\setminus\{v\})$ and $W\cap\operatorname{link}(v)=\varnothing$, so for any $w\in W$ a trace-zero unitary $a_w\in M_w\subseteq M_1$ satisfies $E_B(a_w)=0$ (a single letter from $M_w^o$ with $w\notin\operatorname{link}(v)$ is orthogonal to $L^2B$). We must produce $u_1,u_2\in\mathcal U(M_1)$ with $E_B(u_1)=E_B(u_2)=E_B(u_1^*u_2)=0$. We distinguish cases according to $|W|$; note $W\neq\varnothing$, for otherwise every vertex $\neq v$ is adjacent to $v$, so taking $\GG_1=\{v\}$ and $\GG_2=\GG\setminus\{v\}$ exhibits $\GG$ as reducible, contrary to hypothesis.

\emph{Case A: $|W|\ge2$.} Choose distinct $w_1,w_2\in W$, and set $u_1=a_{w_1}$, $u_2=a_{w_2}$. Then $E_B(u_1)=E_B(u_2)=0$. For $u_1^*u_2=a_{w_1}^*a_{w_2}$: if $(w_1,w_2)\notin\EE(\GG)$ this is the reduced word $a_{w_1}^*a_{w_2}$ of length $2$; if $(w_1,w_2)\in\EE(\GG)$ then $a_{w_1}^*$ and $a_{w_2}$ commute and $u_1^*u_2=a_{w_2}a_{w_1}^*$, again a reduced word using the two letters $w_1,w_2$. In either case every letter of $u_1^*u_2$ lies in $W$, hence outside $\operatorname{link}(v)$, so $u_1^*u_2$ is orthogonal to $L^2B$ and $E_B(u_1^*u_2)=0$.

\emph{Case B: $|W|=1$, say $W=\{w\}$.} Then every vertex $u'\notin\{v,w\}$ is adjacent to $v$ (i.e.\ $u'\in\operatorname{link}(v)$). We claim there is a vertex $u'\notin\{v,w\}$ with $(u',w)\notin\EE(\GG)$. Suppose not: then $w$ is adjacent to every vertex $u'\notin\{v,w\}$. Combined with the fact that every such $u'$ is adjacent to $v$, the partition $\GG_1=\{v,w\}$, $\GG_2=\VV(\GG)\setminus\{v,w\}$ has the property that every $x\in\GG_1$ is adjacent to every $y\in\GG_2$: for $x=v$, $y=u'\in\operatorname{link}(v)$ gives $(v,u')\in\EE(\GG)$; for $x=w$, $(w,u')\in\EE(\GG)$ by assumption. Since $|\VV(\GG)|\ge3$, both parts are nonempty (indeed $|\GG_2|\ge1$), so this exhibits $\GG$ as reducible, contradicting irreducibility. Hence such a $u'$ exists; fix it, and let $b\in M_{u'}$ be a trace-zero unitary. Note $u'\in\operatorname{link}(v)$, so $M_{u'}\subseteq B$.

Set $u_1:=a_w$ and $u_2:=b\,a_w$. Both are unitaries in $M_1$ (since $a_w\in M_w\subseteq M_1$ and $b\in M_{u'}\subseteq M_1$). We verify \eqref{eq:eb0}:
\begin{itemize}
\item $E_B(u_1)=E_B(a_w)=0$, since $w\notin\operatorname{link}(v)$.
\item $E_B(u_2)=E_B(b\,a_w)=b\,E_B(a_w)=0$, since $b\in B$ and $a_w\in M_w^o$ with $w\notin\operatorname{link}(v)$ (the word $b\,a_w$ has reduced form $b\cdot a_w$ with the rightmost letter $a_w\in M_w^o$ outside $B$, hence is $B$-orthogonal).
\item $E_B(u_1^*u_2)=E_B(a_w^*\,b\,a_w)$. Because $(u',w)\notin\EE(\GG)$, the letters $w$ and $u'$ do not commute, so $a_w^*\,b\,a_w$ is a reduced word of length $3$ in the vertices $w,u',w$ (one cannot collapse it: $a_w^*$ and $a_w$ would have to be adjacent across the non-commuting letter $b$, which is impossible). This reduced word contains the letter $w\notin\operatorname{link}(v)$, so it is orthogonal to $L^2B$ and $E_B(u_1^*u_2)=0$.
\end{itemize}
Thus $u_1,u_2$ have the required properties.

\medskip\noindent\textbf{Conclusion.}
In every case the hypotheses of Lemma~\ref{lem:afp} hold for the decomposition \eqref{eq:star-decomp}: $u_1,u_2\in M_1$ and $u_3\in M_2$ satisfy \eqref{eq:eb0}, and $B$ is the range of trace-preserving conditional expectations from $M_1$ and $M_2$. Therefore $M$ is properly proximal relative to $B=\Gamma(\operatorname{link}(v))$. Since $v\in\VV(\GG)$ was arbitrary, $M$ is properly proximal relative to $\Gamma(\operatorname{link}(v))$ for every $v$.

By Lemma~\ref{lem:intersect-pp} applied to the finite family $\{\operatorname{link}(v)\}_{v\in\VV(\GG)}$, $M$ is properly proximal relative to $\Gamma\big(\bigcap_{v\in\VV(\GG)}\operatorname{link}(v)\big)$. Finally $\bigcap_{v\in\VV(\GG)}\operatorname{link}(v)=\varnothing$: a vertex $w$ lying in $\operatorname{link}(v)$ for all $v$ would in particular lie in $\operatorname{link}(w)$, which is impossible since $w\notin\operatorname{link}(w)$. Hence $\Gamma(\varnothing)=\C$, and $M$ is properly proximal relative to $\C$, i.e.\ $M$ is properly proximal. This completes the proof.
\end{proof}

\begin{thebibliography}{DKE24}

\bibitem[BH18]{BH18amenableabsorption}
R\'emi Boutonnet and Cyril Houdayer.
\newblock Amenable absorption in amalgamated free product von {N}eumann algebras.
\newblock {\em Kyoto J. Math.}, 58(3):583--593, 2018.

\bibitem[BIP21]{BIP21}
R\'emi Boutonnet, Adrian Ioana, and Jesse Peterson.
\newblock Properly proximal groups and their von {N}eumann algebras.
\newblock {\em Ann. Sci. \'Ec. Norm. Sup\'er. (4)}, 54(2):445--482, 2021.

\bibitem[CF17]{CaFi17}
Martijn Caspers and Pierre Fima.
\newblock Graph products of operator algebras.
\newblock {\em J. Noncommut. Geom.}, 11(1):367--411, 2017.

\bibitem[DEP23]{DKEP23}
Changying Ding, Srivatsav~Kunnawalkam Elayavalli, and Jesse Peterson.
\newblock Properly proximal von {N}eumann algebras.
\newblock {\em Duke Math. J.}, 172(15):2821--2894, 2023.

\bibitem[DKE24]{DKE24properproximalgroups}
Changying Ding and Srivatsav Kunnawalkam~Elayavalli.
\newblock Proper proximality among various families of groups.
\newblock {\em Groups Geom. Dyn.}, 18(3):921--938, 2024.

\bibitem[Gre90]{Gr90}
Elisabeth~Ruth Green.
\newblock {\em Graph products of groups}.
\newblock PhD thesis, University of Leeds, 1990.

\bibitem[IPR19]{IsPeRu19}
Ishan Ishan, Jesse Peterson, and Lauren Ruth.
\newblock Von {N}eumann equivalence and properly proximal groups.
\newblock arXiv:1910.08682, 2019.

\bibitem[OP10a]{OzPo10a}
Narutaka Ozawa and Sorin Popa.
\newblock On a class of {${\rm II}_1$} factors with at most one {C}artan subalgebra.
\newblock {\em Ann. of Math. (2)}, 172(1):713--749, 2010.

\bibitem[OP10b]{OP10b}
Narutaka Ozawa and Sorin Popa.
\newblock On a class of {${\rm II}_1$} factors with at most one {C}artan subalgebra, {II}.
\newblock {\em Amer. J. Math.}, 132(3):841--866, 2010.

\bibitem[Oza04]{Oza04}
Narutaka Ozawa.
\newblock Solid von {N}eumann algebras.
\newblock {\em Acta Math.}, 192(1):111--117, 2004.

\bibitem[TY25]{toyosawa2025relative}
Kai Toyosawa and Zhiyuan Yang.
\newblock On relative biexactness of amalgamated free product von {N}eumann algebras.
\newblock {\em arXiv preprint arXiv:2505.19508}, 2025.

\end{thebibliography}
\end{document}
